\section{Por qué la robótica también quiere un Foundation Model}

Los campos del lenguaje y la visión ya mostraron la capacidad de reutilización que aporta el entrenamiento a escala; en robótica, en cambio, todavía es común recolectar datos, entrenar y ajustar parámetros por separado para cada tarea. El título inicial y el montage del laboratorio hacen visible esa brecha: los robots logran muchas demostraciones vistosas, pero el crecimiento de sus capacidades sigue pareciendo un conjunto de proyectos aislados entre sí, y no una curva unificada que pueda extenderse de forma sostenida.

\subsection{De un conjunto de Demos robóticas a un problema de Scaling}

La slide de título nombra la dirección como «Towards a Robotics Foundation Model», y en ella `Towards` importa más que `Foundation Model`. El segundo montage muestra la experiencia de varios brazos robóticos en distintas cocinas, objetos y tareas, lo que indica que el control en el mundo real ya cuenta con muchas capacidades locales; el verdadero problema es si un sistema escalable puede absorber en conjunto toda esa experiencia.

\lecturefigure{slide-01-title.jpg}{Towards a Robotics Foundation Model: una agenda de investigación de la clase, no una declaración de algo terminado}{Diapositiva V001 recuperada de la grabación oficial; Stanford CS25 Lecture 15}

\lecturefigure{slide-02-robot-experience-montage.jpg}{La diversidad de la experiencia de robots reales: distintas cocinas, objetos, acciones y embodiment}{Diapositiva V002 recuperada de la grabación oficial; Stanford CS25 Lecture 15}

\teachervoice{Desde el inicio, Ted Xiao describió esta clase como su «intento de entender» las direcciones futuras. Su objetivo es preguntar si las capacidades de los robots pueden pasar de un progreso lineal, por proyectos, a un crecimiento compuesto más rápido, y no afirmar que hoy ya exista una emergence verificable de robots de propósito general.}

\subsection{Ruta de la clase: Ingredient, System, Feedback, Data}

La Agenda primero presenta tres tipos de ingredient y luego pasa, en orden, por RT-1, SayCan, Inner Monologue y DIAL. Este orden corresponde a un robot completo: RT-1 aprende habilidades de bajo nivel, SayCan combina habilidades, Inner Monologue replanifica a partir de la retroalimentación y DIAL amplía la cobertura lingüística de los datos fuera de línea; al final, todas las capas se vuelven a ensamblar en un mapa del sistema.

\lecturefigure{slide-03-agenda-foundation.jpg}{Ruta del curso: de los foundation-model ingredients a cuatro casos de sistemas robóticos}{Diapositiva V003 recuperada de la grabación oficial; Stanford CS25 Lecture 15}

\subsection{El atractivo de la Emergence y la Homogenization}

Los dos puntos más atractivos de un foundation model para la robótica son las \term{emergent capabilities} y la \term{homogenization}. Las primeras se refieren a capacidades complejas que aparecen al aumentar la escala y que no se observan en modelos pequeños; la segunda, a que una misma representación e interfaz de entrenamiento pueda servir a tareas posteriores compuestas. En la slide 4, el lado izquierdo toma prestadas las emergence plots de los modelos de lenguaje, y el derecho usa «habitación—edificio—mundo» para ilustrar la enorme amplitud del espacio de estados de un robot.

\lecturefigure{slide-04-why-foundation-model.jpg}{Por qué se quiere un robotics foundation model: emergence y downstream homogenization}{Diapositiva V004 recuperada de la grabación oficial; Stanford CS25 Lecture 15}

\begin{knowledgebox}{Lectura de la figura: esta no es una scaling curve de robótica}
La curva inferior izquierda proviene de tareas de modelos de lenguaje y sirve para ilustrar que una capacidad puede aparecer después de cierta scale; las elipses de la derecha son solo un diagrama conceptual de la cobertura de entornos robóticos. Juntas plantean una hipótesis: quizá la robótica también necesite superar cierta escala de datos, tareas y modelo para que surja una generalización composicional más fuerte. La clase no presentó un punto crítico propio de la robótica, por lo que esta figura no debe tomarse como una law verificada.
\end{knowledgebox}

\teachervoice{Un estudiante preguntó en ese momento si «ya existe un foundation model para robótica». La respuesta del ponente fue que todavía no está claro, y agregó que quizá el campo de la robótica primero necesite ver su propia emergent-capability curve antes de tener fundamento para llamar foundation model a un sistema.}

\subsection{Por qué la robótica es el último Domino}

El lenguaje, las imágenes, el código y la música pueden obtener enormes corpus estáticos de Internet, mientras que los datos robóticos exigen que dispositivos reales interactúen con el mundo. Cada trayectoria implica tiempo de hardware, costo de operadores, riesgos de seguridad, disposición de objetos y cambios de estado de larga duración. Por ello, la slide del Domino no dice que la robotics no haya avanzado, sino que le falta un mecanismo de producción de datos tan barato y replicable como los token a web-scale.

\lecturefigure{slide-05-why-not-foundation-model.jpg}{Why not: la interacción física convierte a la robotics en la foundation-model domino más difícil de derribar}{Diapositiva V005 recuperada de la grabación oficial; Stanford CS25 Lecture 15}

\begin{warningbox}{«Más datos» no es una recomendación ejecutable}
El valor de un episode robótico depende de lo que aporta de nuevo: una tarea nueva, un objeto nuevo, un fondo nuevo, un embodiment nuevo, un lenguaje nuevo o la repetición de la misma acción. Si solo se agregan trayectorias repetidas en escenarios fáciles, el costo sube sin que necesariamente se amplíe la cobertura de comportamientos; más adelante, las curvas de RT-1 confirmarán directamente la importancia de la task diversity.
\end{warningbox}

\subsection{Descomponer la visión en tres tipos de Ingredient}

En lugar de entrenar directamente un modelo robótico gigantesco, el ponente primero descompone el problema en tres tipos de ingredient operables: la architecture / tokenization tomada del ML scaling; los language y common-sense priors de los Internet-scale models externos; y el paso del aprendizaje robótico desde el costoso ensayo y error en línea hacia grandes offline datasets reutilizables.

\lecturefigure{slide-06-ingredients-overview.jpg}{Los tres tipos de ingredientes de la receta de un robotics foundation model}{Diapositiva V006 recuperada de la grabación oficial; Stanford CS25 Lecture 15}

\begin{importantbox}{Los tres tipos de Ingredient no se sustituyen entre sí}
Una arquitectura de alta capacidad resuelve los problemas de representación e interoperabilidad; los foundation models externos aportan un prior semántico, y los offline robot data aportan acciones reales y dinámica de contacto. Tener solo lenguaje sin capacidad de ejecución, solo trayectorias robóticas sin cobertura semántica o solo un modelo grande sin presupuesto de tiempo real deja al sistema detenido en una sola capa.
\end{importantbox}

\subsection{Resumen de la sección}

La robótica aspira a un foundation model porque la emergence y la homogenization pueden convertir habilidades locales en capacidades reutilizables; sin embargo, el alto costo y la heterogeneidad de los datos físicos impiden copiar tal cual el entrenamiento con datos de Internet. Por eso, la clase descompone la visión en tres ingredientes: arquitectura, prior externo y datos fuera de línea.
